% Generated draft: requires mathematical and source review.
\documentclass[11pt]{article}
\usepackage[margin=25mm]{geometry}
\usepackage{fontspec}
\setmainfont{Latin Modern Roman}
\usepackage{amsmath,amssymb,mathtools,unicode-math}
\setmathfont{Latin Modern Math}
\usepackage{enumitem,needspace,xurl,hyperref}
\setlength{\parindent}{0pt}
\setlength{\parskip}{5pt}
\setlength{\emergencystretch}{3em}
\newcommand{\N}{\mathbb{N}}
\newcommand{\Z}{\mathbb{Z}}
\newcommand{\Q}{\mathbb{Q}}
\newcommand{\R}{\mathbb{R}}
\newcommand{\C}{\mathbb{C}}
\newcommand{\F}{\mathbb{F}}
\newcommand{\A}{\mathbb{A}}
\newcommand{\PP}{\mathbb{P}}
\newcommand{\E}{\mathbb{E}}
\newcommand{\eps}{\varepsilon}
\newcommand{\dd}{\,\mathrm{d}}
\begin{document}
\section*{UnsolvedMath ID: 701782}
\textbf{Source problem code:} catalogue entry
\subsection{Definitions and mathematical statement}
Let \(F\) be a non-abelian free group. The supplied record imposes no restriction on its rank or finite generation. Define its derived subgroup by
\[
[F,F]=\bigl\langle x^{-1}y^{-1}xy : x,y\in F\bigr\rangle.
\]
A \emph{Rota--Baxter operator} on \(F\) is a mapping \(B:F\to F\) satisfying
\[
B(g)B(h)=B\bigl(gB(g)hB(g)^{-1}\bigr)
\qquad\text{for all }g,h\in F.
\]
The question is whether there exists such an operator with
\[
\operatorname{im}(B)=\{B(g):g\in F\}=[F,F].
\]
Here \(F\) is an arbitrary non-abelian free group, as in the supplied formulation; no particular rank is specified.
\subsection{Short English statement}
Given a non-abelian free group \(F\), does there exist a Rota--Baxter operator whose image is exactly the derived subgroup \([F,F]\)? The supplied registry labels this problem open.
\subsection{Sources}
\begin{itemize}
\item Local UnsolvedMath registry, permanent ID \(701782\), source code \texttt{catalogue entry}, title \emph{Kourovka Notebook Problem 20.124}. The supplied record labels the problem open.
\item Attributed in the supplied record to V. G. Bardakov, proposed in 2022.
\item Supplied source attribution: \emph{The Kourovka Notebook (Unsolved Problems in Group Theory)}, 20th issue (2022), problem 20.124; also identified with the 21st issue edition (2026), \texttt{arXiv:1401.0300v48}. The registry describes the collection as edited by E. I. Khukhro and V. D. Mazurov. These bibliographic details are supplied source data and have not been independently verified.
\end{itemize}
\end{document}
